\subsection{Examples}

\subitemhead{Locally contractible spaces}

DEFINITION 3. --- A topological space B is said to be locally contractible if every point b of B has a neighborhood V such that the pointed space (V, b) is contractible (III, p. 234, example 3).

PROPOSITION 5. --- A locally contractible topological space is loop-disentanglable.

Let B be a locally contractible space. Let b be a point of B and let V be a neighborhood of b such that \((V, b)\) is a pointed contractible space. The space V is homotopy equivalent to a point, hence \(\pi_{1}(V, b) = \{\varepsilon_{b}\}\) (III, p. 296, corollary 3).

To prove the proposition, it remains to show that the space B is locally path-connected. Let \(\sigma: V \times I \to V\) be a pointed homotopy at \(b\) relating the constant map with image \(b\) to the map \(Id_V\) . For every neighborhood W of \(b\) contained in V, the set \(\sigma(W \times I)\) is a neighbourhood of \(b\) since it contains \(W = \sigma(W \times \{1\})\) , and it is path-connected because for every \((a, s) \in W \times I\) , the map \(t \mapsto \sigma(a, ts)\) is a path connecting \(b = \sigma(a, 0)\) to \(\sigma(a, s)\) in \(\sigma(W \times I)\) . Let us prove that the sets of the form \(\sigma(W \times I)\) form a fundamental system of neighborhoods of \(b\) . Let \(V'\) be an open neighborhood of \(b\) in B; then, \(\sigma^{-1}(V')\) is an open neighborhood of \(\{b\} \times I\) in \(V \times I\) . Since the space I is compact, the projection \(\mathrm{pr}_1: V \times I \to V\) is proper (TG, I, p. 77, cor. 5) and \(\mathrm{pr}_1(\complement\,\sigma^{-1}(V'))\) is a closed subset of V not containing \(b\) . Its complement W is thus an open neighborhood of \(b\) in V such that \(\mathrm{W} \times \mathbf{I} \subset \sigma^{-1}(\mathrm{V}')\) , whence \(\sigma(\mathrm{W} \times \mathbf{I}) \subset \mathrm{V}'\) .

Every open subset of a numerical space or of a projective space, real or complex, of dimension n (cf. chapters VI and VIII) is loop-disentanglable, as are the Euclidean spheres \(S_{n}\) (TG, VI, p. 11, prop. 4). More generally, every topological manifold (VAR R, second part, p. 7, Notations and Conventions) is loop-disentanglable. Indeed, these spaces are locally contractible.

\subitemhead{Poincaré group of the circle}

The topological space R is homotopic to a point (III, p. 234, example 4), hence simply path-connected (IV, p. 340). The canonical surjection p of R onto T = R/Z makes it a principal covering with group Z (I, p. 100, example 4). The topological space T is loop-disentanglable and (R, 0) is a universal covering of (T, 0). From this we deduce a canonical isomorphism of groups \(\pi_{1}(\mathbf{T}, 0) \to \mathbf{Z}\) : it maps the class of the loop \(t \mapsto p(nt)\) at 0 to the element n of Z. Taking into account example 6 (I, p. 101), we deduce the following proposition.

Proposition 6. — The map \(p \colon x \mapsto e^{2\pi ix}\) from \(\mathbf{R}\) onto \(\mathbf{S}_1\) makes \((\mathbf{R}, 0)\) a universal covering of \((\mathbf{S}_1, 1)\) . The group \(\pi_1(\mathbf{S}_1, 1)\) is isomorphic to \(\mathbf{Z}\) ; the class of the loop \(t \mapsto e^{2\pi it}\) is a generator. For every integer \(n > 0\) , the map \(z \mapsto z^n\) makes \(\mathbf{S}_1\) a Galois covering with group \(\mathbf{Z}/n\mathbf{Z}\) of \(\mathbf{S}_1\) . Every connected and nonempty covering of \(\mathbf{S}_1\) is isomorphic to one of these coverings.

Only the last assertion remains to be proved. The fiber \(E_{1}\) over 1 of a connected and nonempty covering E of \(\mathbf{S}_{1}\) is equipped with a transitive action of the group \(\pi_{1}(\mathbf{S}_{1},1)\) . Since the subgroups of Z are the sets nZ, for \(n \geqslant 0\) , one observes that \(E_{1}\) is isomorphic to one of the homogeneous sets associated with the coverings described previously. Since \(\mathbf{S}_{1}\) is connected and locally path-connected, E is isomorphic to one of these coverings.

COROLLARY. --- The map \(z \mapsto e^{z}\) from C to \(C^{*}\) makes \((\mathbf{C},0)\) a universal covering of \((\mathbf{C}^{*},1)\) . The group \(\pi_{1}(\mathbf{C}^{*},1)\) is isomorphic to Z; the class of the loop \(t \mapsto e^{2\pi it}\) is a generator. For every integer n \textgreater{} 0, the map \(z \mapsto z^{n}\) makes \(C^{*}\) a Galois covering with group Z/nZ of \(C^{*}\) . Every connected and nonempty covering of \(C^{*}\) is isomorphic to one of these coverings.

The map \(z \mapsto (|z|, z/|z|)\) is a homeomorphism of the space \(\mathbf{C}^*\) onto the space \(\mathbf{R}_+^* \times \mathbf{S}^1\) (TG, VI, p. 10, prop. 3); in particular, \(\mathbf{C}^*\) is path-connected. Since the map \(x \mapsto e^x\) from \(\mathbf{R}\) onto \(\mathbf{R}_+^*\) is a homeomorphism, it follows from Proposition 6 that the map \((x, y) \mapsto e^x e^{2\pi iy}\) makes \(\mathbf{R}^2\) a covering of \(\mathbf{C}^*\) . By identifying \(\mathbf{R}^2\) with \(\mathbf{C}\) by the mapping \((x, y) \mapsto x + iy\) , we conclude that the space \(\mathbf{C}\) , equipped with the map \(z \mapsto e^z\) , is a covering of \(\mathbf{C}^*\) . Since the space \(\mathbf{C}\) is connected and simply path-connected, the pointed covering \((\mathbf{C}, 0)\) is a universal covering of the pointed space \((\mathbf{C}^*, 1)\) .

It also follows from the corollary, III, p. 297, that the Poincaré group of \(\mathbf{C}^*\) at 1 is isomorphic to \(\mathbf{Z}\) and that the class \(\gamma\) of the loop \(t \mapsto e^{2\pi it}\) is a generator.

  Let n be an integer \textgreater{} 0. The map \(x \mapsto x^{n}\) from \(R_{+}^{*}\) onto itself is a homeomorphism; we deduce as before that the map \(z \mapsto z^{n}\) makes \(C^{*}\) a covering of degree n of \(C^{*}\) . It is Galois with group Z/nZ. The last assertion is proved as in the proof of Prop. 6.

For every integer \(n \geqslant 0\) , the torus \((\mathbf{S}_{1})^{n}\) is a loop-disentanglable space (IV, p. 341, prop. 1) and its Poincaré group at every point is isomorphic to \(\mathbf{Z}^{n}\) (III, p. 297, prop. 4).

We will generalize these results in Section 3, which is devoted to coverings of topological groups.

\subitemhead{Real projective spaces}

Let n be an integer \textgreater{} 0. The spaces \(S_{n}\) and \(\mathbf{P}_{n}(\mathbf{R})\) are connected, nonempty, and the canonical map (TG, VI, p. 13) \(\varphi\colon\mathbf{S}_{n}\to\mathbf{P}_{n}(\mathbf{R})\) makes \(S_{n}\) a principal covering of \(\mathbf{P}_{n}(\mathbf{R})\) of the group \(\{+1,-1\}\) (I, p. 99, example 1). For \(n\geqslant2\) the sphere \(S_{n}\) is simply connected (I, p. 127, example 3) and loop-disentanglable, hence simply path-connected (IV, p. 344, corollary 1). For n=1, the map \(z\mapsto z^{2}\) from \(S_{1}\) to \(S_{1}\) defines an equivalence relation in \(S_{1}\) whose classes are the pairs of opposite points of \(S_{1}\) . The map \(\varphi\colon\mathbf{S}_{1}\to\mathbf{P}_{1}(\mathbf{R})\) defined by passing to the quotient defines a continuous bijection \(\psi\colon\mathbf{S}_{1}\to\mathbf{P}_{1}(\mathbf{R})\) such that \(\psi(z^{2})=\varphi(z)\) . Since \(S_{1}\) is a compact space, the map \(\psi\) is a homeomorphism (TG, I, p. 63, cor. 2).

PROPOSITION 7. --- The map \(x \mapsto \varphi(e^{2\pi ix})\) from \(\mathbf{R}\) onto \(\mathbf{P}_{1}(\mathbf{R})\) makes \((\mathbf{R},0)\) a universal covering of \((\mathbf{P}_{1}(\mathbf{R}),\varphi(1))\) . Let c: \(\mathbf{I} \to \mathbf{S}_{1}\) be the path \(t \mapsto e^{\pi it}\) ; the class of \(\varphi(c)\) is a generator of the group \(\pi_1(\mathbf{P}_{1}(\mathbf{R}), \varphi(1))\) which is isomorphic to \(\mathbf{Z}\) .

For every integer \(n \geqslant 2\) and every point x of \(S_{n}\) , the canonical map \(\varphi: \mathbf{S}_{n} \to \mathbf{P}_{n}(\mathbf{R})\) makes \((\mathbf{S}_{n}, x)\) a universal covering of \((\mathbf{P}_{n}(\mathbf{R}), \varphi(x))\) . For every path c in \(S_{n}\) joining x to -x, the class of \(\varphi \circ c\) generates the group \(\pi_{1}(\mathbf{P}_{n}(\mathbf{R}), \varphi(x))\) which is isomorphic to Z/2Z.

\subitemhead{Quotient of a space by a proper action of a discrete group}

Lemma 1. --- Let X be a topological space, let G be a discrete group acting properly on X, and let p be the canonical projection from X onto X/G.

Let \(x\) be a point of \(X\) , let \(K_x\) be its fixator in \(G\) . Let \(U_1\) be a neighborhood of \(x\) in \(X\) ; there exists a neighborhood \(U\) of \(x\) in \(X\) stable under \(K_x\) and contained in \(U_1\) such that \(g \cdot U \cap U = \varnothing\) if \(g \notin K_x\) and such that the canonical map \(p\) induces a homeomorphism of \(U / K_x\) onto a neighborhood \(V\) of \(p(x)\) in \(X / G\) .

By prop. 8 of TG, III, p. 32, there exists a neighborhood \(U_{2}\) of x in X, stable under \(K_{x}\) , such that \(g \cdot U_{2} \cap U_{2} = \varnothing\) if \(g \notin K_{x}\) and such that the canonical map p induces a homeomorphism from \(U_{2}/K_{x}\) onto a neighborhood of \(p(x)\) in X/G.

As the canonical map \(U_{2} \rightarrow U_{2}/K_{x}\) is closed (TG, III, p. 28, prop. 2), there exists an open neighborhood U of x contained in \(U_{1} \cap U_{2}\) which is stable under \(K_{x}\) . The equivalence relation in \(U_{2}\) defined by \(K_{x}\) is also open (TG, III, p. 10, lemma 2). It then follows from TG, I, p. 32, prop. 4, that the canonical mapping induces a homeomorphism from \(U/K_{x}\) onto an open neighborhood of \(p(x)\) in X/G, whence the lemma.

PROPOSITION 8. --- Let X be a topological space and let G be a discrete group acting properly on X. Denote by p the canonical projection from X to X/G.

\begin{enumerate}
\def\labelenumi{\alph{enumi})}
\item
  Suppose that X is path-connected and that the group G is generated by the stabilizers of the points of X. Then the canonical homomorphism \(\pi_{1}(X,x)\to\pi_{1}(X/G,p(x))\) is surjective for every point x in X. In particular, X/G is simply path-connected if X is.
\item
  If the space X is loop-disentanglable, then the space X/G is loop-disentanglable.
\item
  The set N of elements \(g \in G\) for which there exists a path \(c_{g} \colon I \to X\) such that \(c_{g}(0) = x\) , \(c_{g}(1) = g \cdot x\) and such that the loop \(p \circ c_{g}\) in X/G is strictly homotopic to a constant loop is a subgroup of G. If \(g \in G\) and if there exists a point y of X such that \(g \cdot y = y\) , choose a path c: I → X joining x to y; then the path \(c' = c * (g \cdot c)\) satisfies \(c'(0) = x\) , \(c'(1) = g \cdot x\) and \([p \circ c'] = [p \circ c][p \circ (g \cdot c)]^{-1} = e_{p(x)}\) , therefore \(p \circ c'\) is strictly homotopic to a constant loop. Since G is generated by the fixators of the points of X, it follows that N = G.
\end{enumerate}

Let \(c\) be a loop in X/G at \(p(x)\) . By Theorem 4 of III, p. 287, there exists a path \(\widetilde{c} \colon \mathbf{I} \to \mathrm{X}\) lifting \(c\) such that \(\widetilde{c}(0) = x\) . Since \(p(\widetilde{c}(1)) = c(1) = p(x)\) , there exists \(g \in \mathrm{G}\) such that \(\widetilde{c}(1) = g \cdot x\) . Choose a path \(c_g \colon \mathbf{I} \to \mathrm{X}\) connecting \(x\) to \(g \cdot x\) and such that \(p \circ c_g\) is strictly homotopic to a constant loop. Then the path \(c' = \widetilde{c} * \overline{c_g}\) is a loop at \(x\) in X such that \([p \circ c'] = [c]\) . This shows that the homomorphism \(\pi_1(\mathrm{X}, x) \to \pi_1(\mathrm{X/G}, p(x))\) is surjective. The other assertion follows immediately.

Let us now prove assertion \(b\) ). The space X/G is locally path-connected (III, p. 261, prop. 8).

Let \(x\) be a point of X and let \(K_x\) be its stabilizer in G. Let \(U_1\) be an open neighborhood of \(x\) such that the image of the canonical homomorphism \(\pi_1(U_1, x) \to \pi_1(X, x)\) is reduced to the identity element. By lemma 1 (IV, p. 349), there exists a neighborhood U of \(x\) in X, contained in \(U_1\), stable under \(K_x\) , such that \(g \cdot U \cap U = \varnothing\) if \(g \notin K_x\) and such that the canonical map \(p\) induces a homeomorphism of \(U/K_x\) onto a neighborhood V of \(p(x)\) in X/G.

Let \(c\) be a loop at \(p(x)\) in V; by theorem 4 (III, p. 287), applied to the topological space U and the group \(K_x\) , there exists a path \(\widetilde{c} \colon \mathbf{I} \to \mathbf{U}\) such that \(\widetilde{c}(0) = x\) and which lifts \(c\) . One necessarily has \(\widetilde{c}(1) = x\) , so that \(\widetilde{c}\) is a loop at \(x\) in U. By hypothesis, \(\widetilde{c}\) is strictly homotopic to a constant loop in X; consequently, the same holds for \(c\) and the canonical homomorphism \(\pi_1(V, p(x)) \to \pi_1(X/G, p(x))\) is trivial. This shows that X/G is loop-disentanglable if X is.

\begin{enumerate}
\def\labelenumi{\arabic{enumi})}
\setcounter{enumi}{4}
\tightlist
\item
  In the Euclidean plane \(R^{2}\) , let P be the topological space which is the union of the circles with center \((1/n,0)\) passing through the origin (for \(n \in N^{*}\) ). The space P is compact, connected, and locally path-connected, but is not loop-disentanglable. The group \(\pi_{1}(P,0)\) , equipped with the admissible topology, is Hausdorff and non-discrete (III, p. 337, exerc. 7).
\end{enumerate}

Remark 1. --- Theorem 4 of III, p. 287 and prop. 8 (IV, p. 349) remain valid under the more general hypothesis that G is a finite-dimensional Lie group over R acting properly on X. For more details, cf. D. MONTGOMERY and C. T. YANG, “The existence of a slice”, Annals of mathematics 65 (1957), p. 108--116; R. PALAIS, “On the existence of slices for actions of non-compact Lie groups”, Annals of mathematics 73 (1961), p. 295--323; and G. BREDON, Introduction to compact transformation groups, Academic Press, 1972.

